\documentclass[10pt,a4paper]{amsart}
\usepackage[margin=19mm]{geometry}
\usepackage{amsmath,amssymb,mathtools}
\usepackage[scr=rsfs,cal=euler]{mathalfa}
\usepackage{xfrac}
\usepackage[unicode,hidelinks,bookmarks=false]{hyperref}
\newcommand{\ie}{i.e.\ }
\newcommand{\cf}{cf.\ }
\newcommand{\resp}{respectively\ }
\newcommand{\Subsection}{\subsection}
\providecommand{\nobreakdash}{\nobreak\hbox{-}}
\setlength{\emergencystretch}{2em}
\multlinegap0pt
% ---------------------------------------------------------------------------
% Editorial AMS wrapper prelude. The theorem environments and the mathematical macros below are
% transcribed from the paper's own preamble (cached-originals/source0.tex lines 28--62); the AMS amsart
% class, the named format packages of the pinned TeX Live runtime and this editorial formatting are not
% part of the author's text. No mathematics is changed.
% ---------------------------------------------------------------------------
\newtheorem{thm}{Theorem}[section]
\newtheorem{cor}[thm]{Corollary}
\newtheorem{lem}[thm]{Lemma}
\newtheorem{prop}[thm]{Proposition}
\newtheorem{rem}[thm]{Remark}
\newtheorem{defi}[thm]{Definition}
\newtheorem{exa}[thm]{Example}
\newcommand{\dom}{\mathrm{dom}}
\newcommand{\rng}{\mathrm{rng}}
\newcommand{\N}{\mathbb N}
\newcommand{\Z}{\mathbb Z}
\newcommand{\Q}{\mathbb Q}
\newcommand{\R}{\mathbb R}
\newcommand{\C}{\mathbb C}
\newcommand{\Bai}{\N^{\N}}
\newcommand{\Seq}{\N^{<\omega}}
\newcommand{\Concil}{\N^{\leq\omega}}
\newcommand{\Can}{2^{\N}}
\newcommand{\seq}{2^{<\omega}}
\newcommand{\concil}{2^{\leq\omega}}
\newcommand{\bsigma}{\boldsymbol\Sigma}
\newcommand{\bpi}{\boldsymbol\Pi}
\newcommand{\bdelta}{\boldsymbol\Delta}
\newcommand{\conc}{{}^\smallfrown}
\newcommand{\lnh}{{\rm{length}}}
\newcommand{\card}{{\rm{card}}}
\newcommand{\wadge}{\le_W}

\newcommand{\quot}[2]{{\raisebox{.2em}{$#1\!$}\left/\raisebox{-.2em}{$#2$}\right.}}

\newcommand{\Cl}[1]{\mathsf{Cl}(#1)}
\newcommand{\Int}[1]{\mathsf{Int}(#1)}
\newcommand{\Fr}[1]{\mathsf{Fr}(#1)} 





\def\eg{{\em e.g.}}
\def\ie{{\em i.e.}}
\def\cf{{\em cf.}}

\begin{document}
\title{The complexity of being monitorable}
\author{Riccardo Camerlo and Francesco Dagnino}
\date{Source: arXiv:2601.04256v3; distributed under CC BY 4.0}
\maketitle
\noindent\emph{Source, version and licence.} \emph{The complexity of being monitorable}, by Riccardo Camerlo and Francesco Dagnino. The source of this editable unit is the e-print of \textbf{version v3} of arXiv:2601.04256, https://arxiv.org/abs/2601.04256v3, whose material is distributed under the Creative Commons Attribution 4.0 International licence, http://creativecommons.org/licenses/by/4.0/, as recorded in the arXiv licence metadata for this paper and shown on that abstract page. The whole of the body below is version v3, the newest version of the paper. Full rights notice: licenses/arxiv-2601.04256-CC-BY-4.0.txt.\par\smallskip
\noindent\textit{Editorial scope.} Counted unit: original Theorem 3.18 of the article, the statement carried by the source environment \texttt{thm} with source label \texttt{thm:trn}, cached-originals/source0.tex lines 726--731, printed as \textbf{Theorem 3.18} exactly as in the article. It is delivered together with the authors' own complete proof as the single primary proof body, source0.tex lines 733--760 (28 lines): the \texttt{proof} environment that immediately follows the statement, proves item (1) through the Borel description of $\mathcal R_1$ and item (2) through the dense $\bpi^0_2$ set of relations with hyperconnected tree, and closes with the density argument. No proof marker is added and no mathematics is written, repaired or re-derived; the authors' notation and the printed slips of the source are kept literal. Necessary complete context is retained uncounted, in source order: source0.tex line 366, the subsection heading of the second countable case that the counted proof refers to for its characterisation; lines 580--582, the lemma producing the integer $N$ used by that characterisation, with its complete proof at lines 584--599; and lines 604--606, Lemma 3.13 on hyperconnected spaces, which identifies the dense set of part (2). The article's own numbering is reproduced by explicit counter settings: the article's publisher class declares \texttt{\textbackslash newtheorem\{thm\}\{Theorem\}[section]} with corollary, lemma, proposition, remark, definition and example sharing that counter, so the wrapper restates those declarations and sets the counter before each retained passage; the counted statement prints as Theorem 3.18 and the cited lemma as Lemma 3.13, exactly as in the article. Every \texttt{\textbackslash ref} and \texttt{\textbackslash eqref} inside every retained passage resolves inside this delivered file. No \texttt{\textbackslash cite} command occurs in any retained passage: the closure of the retained passages over references and citations was computed, it contains no reference and no citation at all, so this unit delivers no bibliography entry and the article's own bibliography data file \texttt{biblio.bib} is neither spliced into the bundled source nor redistributed. The retained passages contain no figure environment and no graphics-inclusion command, so no artwork is redistributed. The source's own class is the publisher class \texttt{lmcs} of Logical Methods in Computer Science (\texttt{\textbackslash documentclass\{lmcs\}}); that class file is shipped by the e-print but is \emph{not} shipped, loaded or used by this package and is never redistributed, and the wrapper is the AMS \texttt{amsart} class with the article's theorem declarations restated and the article's own macros transcribed from its preamble. Every declared body of this unit is an exact, unmodified span of source0.tex: no body is a bounded passage comparison, \texttt{omit} was not used anywhere, and consequently no fidelity receipt is asserted in delivery-evidence.json. The complete concatenated original source is bundled as sources/192285/source0.tex. Omitted from this editable unit and therefore absent from it are source0.tex lines 1--365; 367--579; 583--583; 600--603; 607--725; 732--732; 761--850, that is the article's own preamble, title and author block and abstract, the introduction, the remaining propositions, lemmas and corollaries with their proofs, the non-second-countable case, and the article's bibliography data, all of which remain present in the bundled source but are not delivered here. The AMS \texttt{amsart} wrapper, the named format packages of the pinned TeX Live runtime and this editorial source, licence and scope paragraph are editorial formatting only.\par\smallskip
\setcounter{section}{3}
\setcounter{equation}{0}
\setcounter{thm}{0}
\subsection{The second countable case} \label{sec:sccnt}

\setcounter{section}{3}
\setcounter{equation}{1}
\setcounter{thm}{11}
\begin{lem}
There exists $N\in \N $ such that every family of pairwise disjoint open subsets of $L$ has cardinality at most $N$.
\end{lem}

\setcounter{section}{3}
\setcounter{equation}{1}
\setcounter{thm}{12}
\begin{proof}
Note that if $ \mathcal W $ is a family of non-empty, pairwise disjoint, open subsets of a topological space, then $ \mathcal W $ is maximal if and only if $\bigcup \mathcal W $ is dense.

Assume towards contradiction that there exist families of arbitrarily high finite cardinality of pairwise disjoint open subsets of $L$.
We define by recursion a sequence of maximal families $ \mathcal W_n$ of non-empty, pairwise disjoint, open subsets of $L$, with each $ \mathcal W_{n+1}$ refining $ \mathcal W_n$.

Let $ \mathcal W_0=\{ L\} $.

Assume $ \mathcal W_n$ is defined, and let $ \mathcal V $ be a maximal family of non-empty, pairwise disjoint, open subsets of $L$, with $ \card ( \mathcal V)> \card ( \mathcal W_n)$.
Set $ \mathcal W_{n+1}=\{ W\cap V\mid W\in \mathcal W_n,V\in \mathcal V ,W\cap V\ne\emptyset\} $; notice that $ \mathcal W_{n+1}$ is maximal and $ \card ( \mathcal W_{n+1})\geqslant \card ( \mathcal V )> \card ( \mathcal W_n)$.

Let $ \mathcal T $ be the tree whose nodes are the elements of $\bigcup_{n\in \N } \mathcal W_n$ with the order of reverse inclusion.
Since there are infinitely many branching nodes and $ \mathcal T $ is finite splitting, let $b$ be a branch of $ \mathcal T $ containing infinitely many of them.
If $W$ is a branching point belonging to $b$, let $W^+$ be an immediate successor of $W$ that does not belong to $b$.
Then the family of all such $W^+$ is infinite and consists of pairwise disjoint open subsets of $L$, a contradiction.
\end{proof}

\setcounter{section}{3}
\setcounter{equation}{1}
\setcounter{thm}{12}
\begin{lem} \label{lem:newcorr}
If $Y$ is a countable, second countable, hyperconnected space, then $ \mathcal M (Y)\in \bsigma^0_2( \mathcal P (Y))$.
\end{lem}

\setcounter{section}{3}
\setcounter{equation}{1}
\setcounter{thm}{17}
\begin{thm} \label{thm:trn}
\begin{enumerate}
\item\label{thm:trn:1} Both $ \mathcal R_0, \mathcal R_1$ are Borel subsets of $ \mathcal P (X\times E\times X)$
\item\label{thm:trn:2} The set $ \mathcal R_0$ is comeagre in $ \mathcal P (X\times E\times X)$; equivalently, the set $ \mathcal R_1$ is meagre in $ \mathcal P (X\times E\times X)$
\end{enumerate}
\end{thm}

\setcounter{section}{3}
\setcounter{equation}{1}
\setcounter{thm}{18}
\begin{proof}
(\ref{thm:trn:1})  Given $q\in X$ and a finite subset $F$ of $E^{<\omega }$, the set $\{ R\in \mathcal P (X\times E\times X)\mid q\in U_F^R\} $ is open.
To see that $ \mathcal R_1$ is Borel, notice that, by the characterisation of section \ref{sec:sccnt}, one has $R\in \mathcal R_1$ if and only if for every $N\in \N $ there exist finite subsets $F_0,\ldots ,F_N$ of $E^{<\omega }$ such that
\begin{itemize}
\item every $U_{F_h}^R$ is non-empty
\item if $h\ne k$ then $U_{F_h}^R\cap U_{F_k}^R=\emptyset $
\item no $U_{F_h}^R$ contains a minimal non-empty, open set; this means that for every $F$ finite subset of $E^{<\omega }$ the following implication holds:
\[
\emptyset\ne U_F^R\subseteq U_{F_h}^R\Rightarrow\exists G\ \emptyset\ne U_G^R\subset U_F^R
\]
\end{itemize}

As the listed conditions are Borel, $ \mathcal R_1$ is Borel, whence $ \mathcal R_0$ is Borel as well.

(\ref{thm:trn:2}) To prove that $ \mathcal R_0$ is comeagre in $ \mathcal P (X\times E\times X)$ we show that it contains a dense $ \bpi^0_2$ set, namely $ \mathcal H =\{ R\in \mathcal P (X\times E\times X)\mid \mathcal T_R \text{ is hyperconnected} \} $ (see lemma \ref{lem:newcorr}).

Note that $R\in \mathcal H $ if and only if, for every $F,G$ finite subsets of $E^{<\omega }$, there exists $q\in U_F^R\cap U_G^R=U_{F\cup G}^R$.
Since the condition $q\in U_{F\cup G}^R$ is open in the variable $R$, it follows that $ \mathcal H \in \bpi^0_2( \mathcal P (X\times E\times X))$.

To show that $ \mathcal H $ is dense, fix a basic open set
\[ \begin{array}{l}
\mathcal U =\{ R\in \mathcal P (X\times E\times X)\mid \\
(q_1,e_1,q'_1),\ldots ,(q_n,e_n,q'_n)\in R, (q''_1,e'_1,q'''_1),\ldots ,(q''_m,e'_m,q'''_m)\notin R\}
\end{array} \]
Let $q\in X\setminus\{ q''_1\ldots ,q''_m,q'''_1,\ldots ,q'''_m\} $ and set $R=\{ (q_1,e_1,q'_1),\ldots ,(q_n,e_n,q'_n)\}\cup\{ (q,e,q)\mid e\in E\} $.
Then $R\in \mathcal U $.
Moreover, $R\in \mathcal H$ since $q\in U_F^R$ for every finite subset $F$ of $E^{<\omega }$.
\end{proof}

\end{document}
